\noindent\hspace*{1.2em}\kbvlabel{ಉದಾ}\hspace{0.6em}ಜೂಜುಗಾರ, ಕಲಾಕಾರ, ಕನ್ನಡಿಗ, ಹಣವಂತ, ಮಾಂತ್ರಿಕ, ಮಡಿವಾಳ, ಚಮ್ಮಾರ, ಹಾವಾಡಿಗ ಇತ್ಯಾದಿ ; ಇವು ಪುಲ್ಲಿಂಗ ತದ್ಧಿತಾಂತ ನಾಮಗಳು.\par

\vspace{2.5pt plus 3pt minus 0.5pt}\noindent\hspace*{1.2em}\kbvnum{೧.}\hspace{0.55em}ಸ್ತ್ರೀಲಿಂಗದಲ್ಲಿ ಗಾತಿ, ವಂತೆ, ಗಿತ್ತಿ, ಇತಿ, ವತಿ, ಎ-ಇತ್ಯಾದಿ ತದ್ಧಿತ ಪ್ರತ್ಯಯಗಳು ನಾಮಪದಗಳ ಮೇಲೆ ಬಂದು ಸ್ತ್ರೀಲಿಂಗ ತದ್ಧಿತಾಂತ ನಾಮಗಳಾಗುವುವು. ಉದಾಹರಣೆಗೆ ಮೋಜುಗಾತಿ, ದಯಾವಂತೆ, ಹೂವಾಡಿಗಿತ್ತಿ, ಕನ್ನಡಿತಿ, ಗುಣವತಿ, ಜಾಣೆ ಇತ್ಯಾದಿ.\par

\kbvsechead{(೨)\hspace{0.45em}ತದ್ಧಿತಾಂತ ಭಾವನಾಮ :}

\indent ಷಷ್ಠೀ ವಿಭಕ್ತಿಗಳಾದ ನಾಮಪದಗಳ ಮೇಲೆ ಭಾವಾರ್ಥದಲ್ಲಿ ತನ, ಇಕೆ, ಮೆ, ಪು, ಉ ಇತ್ಯಾದಿ ತದ್ಧಿತ ಪ್ರತ್ಯಯಗಳು ಸೇರಿ ಆಗುವ ಹೊಸ ಶಬ್ದಗಳಿಗೆ ತದ್ಧಿತಾಂತ ಭಾವನಾಮಗಳೆಂದು ಹೆಸರು. ಉದಾ : ದಡ್ಡ ತನ.\par

\begin{kbvnote}
\noindent{\kbvsecfont ವಿವರಣೆ :}\hspace{0.5em}“ದಡ್ಡ” ಎಂಬ ನಾಮಪದದ ಷಷ್ಠೀ ವಿಭಕ್ತಿರೂಪ “ದಡ್ಡನ” ; ದಡ್ಡನ ಭಾವವೇ ದಡ್ಡ ತನ. ಇಲ್ಲಿ “ದಡ್ಡನ” ಎಂಬ ನಾಮಪದದ ಮುಂದೆ ಭಾವಾರ್ಥದಲ್ಲಿ “ತನ” ಎಂಬ ಪ್ರತ್ಯಯ ಸೇರಿ “ದಡ್ಡ ತನ” ಆಗಿರು ವುದು ; ಇದು ತದ್ಧಿತಾಂತ ಭಾವನಾಮ.
\end{kbvnote}

\indent ಇದೇ ರೀತಿಯಾಗಿ ಗೌಡಿಕೆ, ಜಾಣ್ಮೆ, ತಂಪು, ಕುರುಡು ಇವು\par

\kbvsechead{ತದ್ಧಿತಾಂತ ಭಾವನಾಮಗಳು :}

\kbvsechead{ವಿಶ್ಲೇಷಣೆ :}

\indent - ಗೌಡನ ಭಾವ ಗೌಡಿಕೆ (ಇಲ್ಲಿ “ಇಕೆ” ಪ್ರತ್ಯಯ ಸೇರಿದೆ).\par

\indent - ಜಾಣನ ಭಾವ ಜಾಣ್ಮೆ (ಇಲ್ಲಿ “ಮೆ” ಪ್ರತ್ಯಯ ಸೇರಿದೆ).\par

\indent - ತಣ್ಣನೆಯದರ ಭಾವ ತಂಪು (ಇಲ್ಲಿ “ಪು” ಪ್ರತ್ಯಯ ಸೇರಿದೆ).\par

\indent - ಕುರುಡನ ಭಾವ ಕುರುಡು (ಇಲ್ಲಿ “ಉ” ಪ್ರತ್ಯಯ ಸೇರಿದೆ).\par

\vspace{2.5pt plus 3pt minus 0.5pt}\noindent\hspace*{1.2em}\kbvnum{(೩)}\hspace{0.55em}ತದ್ಧಿತಾಂತಾವ್ಯಯ\par

\begin{kbvderiv}
\vspace{4pt plus 3pt minus 1pt}
\begin{center}\begin{adjustbox}{max width=\textwidth}
{\renewcommand{\arraystretch}{1.22}\setlength{\tabcolsep}{4pt}
\begin{tabular}{l c l}
ಶಂಕರನ + ಅಂತೆ & {\color{kbvgrey}$=$} & \kbvresult{ಶಂಕರನಂತೆ}{\color{kbvgrey} (ಪು.ಎ.)} \\
ಗೌರಿಯ + ಅಂತೆ & {\color{kbvgrey}$=$} & \kbvresult{ಗೌರಿಯಂತೆ}{\color{kbvgrey} (ಸ್ತ್ರೀ.ಎ.)} \\
ಲತೆಯ + ಅಂತೆ & {\color{kbvgrey}$=$} & \kbvresult{ಲತೆಯಂತೆ}{\color{kbvgrey} (ನ.ಎ.)} \\
ಮಕ್ಕಳ + ಅಂತೆ & {\color{kbvgrey}$=$} & \kbvresult{ಮಕ್ಕಳಂತೆ}{\color{kbvgrey} (ನ.ಬ.)} \\
\end{tabular}}
\end{adjustbox}\end{center}
\vspace{4pt plus 3pt minus 1pt}
\end{kbvderiv}
